%% language-test — multilingual paper (acmart `language` option).
%% Main language French (last in the babel list, after the auto-seeded english),
%% with an English secondary title, abstract and keywords via \translatedtitle /
%% translatedabstract / \translatedkeywords. Exercises the localized fixed
%% strings — \keywordsname ("Mots Clés et Phrases Supplémentaires"), \proofname
%% ("Démonstration"), \acksname ("Remerciements") — plus French hyphenation.
%% Twin of language-test.typ.
\documentclass[acmsmall,language=french]{acmart}

\acmJournal{JACM}
\acmVolume{37}
\acmNumber{4}
\acmArticle{111}
\acmYear{2018}
\acmMonth{8}
\acmDOI{XXXXXXX.XXXXXXX}
\setcopyright{acmlicensed}
\copyrightyear{2018}

\begin{document}

\title{Une note sur la complexité de calcul}
\translatedtitle{english}{A note on computational complexity}

\author{Jean Dupont}
\email{dupont@example.fr}
\affiliation{%
  \institution{Institut de Recherche}
  \city{Paris}
  \country{France}}

\renewcommand{\shortauthors}{Dupont}

\begin{abstract}
  Nous présentons une courte note en français afin d'exercer le mode multilingue
  de la classe acmart, avec des chaînes de caractères traduites.
\end{abstract}

\begin{translatedabstract}{english}
  We present a short note, in English, to exercise the multilingual mode of the
  acmart class with translated fixed strings.
\end{translatedabstract}

\keywords{complexité, algorithmes, calcul}
\translatedkeywords{english}{complexity, algorithms, computation}

\maketitle

\section{Introduction}
Ce document vérifie que les chaînes localisées et le contenu traduit du haut de
page correspondent entre la classe LaTeX et le portage Typst.

\begin{proof}
Voici une démonstration qui se termine par un carré CQFD.
\end{proof}

\begin{acks}
Nous remercions les relecteurs.
\end{acks}

\end{document}
